
\subsubsection{Strong Chain Rule}\label{sec:strong-chain-rule}
\iffalse

\begin{itemize}
\item It can be computed element-by-element, starting from the last
element of the distribution.

\item Conditioned on the spoiling leakage, we can simulate the
adversary's entire view in the protocol
(by leaking extra information along with the spoiled bits), other than the outcome of the
random variable $r$ over choice of $R$, where $R$ is drawn from the
probability distribution that conditions on the spoiling leakage.

\item Conditioned on the spoiling leakage, each element $R_i$ either has
max-entropy at most $\sim 1.5 \cdot \lg(n)$ or min-entropy at least
$\sim 1.5 \cdot \lg (n)$.
\end{itemize}

We further show that, due to the high min-entropy of the entire set $R$,
there must be a constant fraction of blocks $R'$ with min-entropy at
least $1.5 \cdot \lg (n)$. To reflect the fact that only a constant
fraction of elements are guaranteed to have high min-entropy, the
parameter $n'$ is set to $n/3$.  We now consider the marginal
distribution over $R'$ (the elements of $R$ with high min-entropy) and
show that the geometric collision property holds with respect to this
distribution.
\fi

Our strong chain rule considers min-entropy where leakage functions
$\ell_1(\cdot), \ldots, \ell_n(\cdot)$ are additionally considered. We describe
our theorem in a general way, and we hope that it may find future applications
in leakage-resilient cryptography.

\paragraph{Sequence of random variables.}
Recall that $R$ is the set of secret items in the min-hash protocol. Here, we
treat $R$ as a sequence of block-by-block random variables $R = (R_1,\dots,R_n)$, associated with 
(potentially randomized) leakage functions $\ell_1(\cdot), \ldots,
\ell_n(\cdot)$ with randomness $\rho_1, \ldots, \rho_n$. %
You can think of the blocks as coming in a streaming fashion in order of $R_1, R_{2}, \ldots, R_n$. 

\paragraph{Leakage functions.}
Loosely speaking, the properties we require of the leakage functions are that
the $i$-th leakage $\ell_i$ can be computed given $(R_i, \rho_i)$, and all the
outputs of $(\ell_{i+1}, \ell_{i+2}, \ldots, \ell_n)$. We also require that the
total number of valid sequences of leakages from $\ell_1(\cdot), \ldots,
\ell_n(\cdot)$ should be sufficiently small (see Property 1 in
Theorem~\ref{thm:chain_leak_3} below).

\paragraph{Spoiling functions.}
Our theorem below states the existence of a spoiling function $f(\cdot)$ with
certain properties, as well as properties of the random variables
$(R_1,\dots,R_n)$ and $(\rho_1, \ldots, \rho_n)$ conditioned on the output of
the spoiling function $f(R)$.

The properties of $(R_1,\dots,R_n)$ and $(\rho_1, \ldots, \rho_n)$ are roughly the following:
(1) There exist disjoint sets $V, W$ such that $V \cup W = [n]$ that are determined by $f(R)$.
(2) Blocks $\{R_i\}_{i \in V}$ have high min-entropy conditioned on $f(R)$. 
(3) Blocks $\{R_i\}_{i \in W}$ have small support size (low max-entropy) conditioned on $f(R)$.
(4) For $i \in V$, the random strings
$\rho_i$ are uniform random and independent conditioned on $f(R)$.
(See Properties (5)-(8) in Theorem~\ref{thm:chain_leak_3}).

The properties of $f(\cdot)$ are roughly the following:
(1) The failure probability (outputting $\bot$) is small.
(2) As long as the total number of valid sequences of leakages
from $\ell_1(\cdot), \ldots, \ell_n(\cdot)$ is sufficiently small, the image size of $f$ is small.
This property ensures that we do not lose too much of the total min-entropy of $R$ by releasing $f(R)$
(3) The leakages $\{\ell_i(\cdot)\}_{i \in W}$ can be computed given $f(R)$.
(See Properties (2)-(4) in Theorem~\ref{thm:chain_leak_3})

\paragraph{Our theorem.}
The main difference between our spoiling lemma and prior ones is that our min and max entropy guarantees on $R = (R_1,\dots,R_n) \mid f(R)$
hold even with respect to additional leakage $\{\ell_i\}_{i \in W}$ which is included in the spoiled bits $f(R)$.

\begin{theorem} [Block structures with few bits spoiled and leakage]\label{thm:chain_leak_3}  
Let $\mathcal{U} = U_1 \times \cdots \times U_n$ be a fixed universe
and $R = (R_1,\dots,R_n)$ be a sequence of (possibly correlated) random variables where each $R_i$ is over $U_i$ (and all are disjoint) and $|U_i| = \ell$ for all $i$.
Let $\rho_1, \ldots, \rho_n$ be a sequence of uniformly
random strings over $\{0,1\}^m$
and
let $\ell_1(\cdot), \ldots, \ell_n(\cdot)$ be leakage functions.
%whose inputs are defined recursively as follows:
%$\ell_n$ receives input $R_n, \rho_n$ and outputs $\beta_n$.
%$\ell_i$ for $i < n$ receives input $R_i, \rho_i$,
%and $\beta_j$, $i < j \leq n$
%and outputs $\beta_i$.
Then, for any $\epsilon \in (0, 1)$, 
any $\dd > 0$ and any $c \in [2^\delta, \ell/2^\delta]$, 
there exists a spoiling 
leakage function $f(R)$ that satisfies the following properties. 
\begin{enumerate}
\iffalse
\item %We will omit the subscript of $f$ for brevity. 
Let $Im(f)$ be the set of images of $f$ and $Dom(f)$ be the domain of $f$.  
Images except $\bot$ have at least $\frac{\epsilon |Dom(f)|}{|Im(f)|}$ pre-images.   \mingyu{Domain, why do we need to lower bound the pre-images?}
\dnote{This takes care of the partitions with low weight. We ultimately need this, along with $|Im(f)|$ to lower bound the remaining min-entropy in $X$ after revealing $y$.}
\fi
    \item A sequence $\beta_1, \ldots, \beta_n$ is valid if 
    for all $i \in V$, $\beta_i = \bot$ and 
    for all $i \in W$, $\beta_i = \ell_i(R_i, \rho_i, \beta_{>i})$, where $\beta_{>i} = (\beta_{i+1}, \ldots, \beta_n)$. 
    We require that the number of valid sequences
    $\beta_1, \ldots, \beta_n$ is at most $B$.
    \item It holds that $\Pr_{R}[f(R) = \bot] \leq \epsilon n$.
    \item $|Im(f)| \leq  B \cdot \left ( 2(\lg(\ell) + \lg (1/\epsilon)) /\delta \right )^{n}$.
    \item Conditioned on any $y \in Im(f) \setminus \{\bot\}$, for all $i \in W$, the leakage 
    $\ell_i(R_i, \rho_i, \beta_{>i})$  
    can be computed from $y$. Here,  $\beta_j = \bot$ if $j \in V$
    and $\beta_j = \ell_j(R_j, \rho_j, \beta_{> j})$ otherwise.
    \item 
    Let $Im(f)$ be the set of images of $f$. 
    Every $y \in Im(f) \setminus \{\bot\}$ specifies two disjoint sets $V$ and $W$ such that $V \cup W = [n]$.
    \item 
    Conditioned on any $y \in Im(f) \setminus \{\bot\}$, 
    for every $i \in V$, every element in distribution $R_i \mid R_{<i}$ has low probability weight, i.e.,
    \begin{align*}
    \forall y \in Im(f) \setminus \{\bot\}, \forall r \mbox{ s.t. } f(r) = y, \forall i \in V:~~ 
    \Pr\left[R_i = r_i~ \Bigg |~R_{<i} = r_{<i},~ y\right]
    \leq  \frac{2^{\delta}}{c}.
    \end{align*}
    \item Conditioned on any $y \in Im(f) \setminus \{\bot\}$, 
    for every $i \in W$, it holds that $R_i \mid R_{<i}$ has small support size, i.e.,
    \begin{eqnarray*}
        \forall y \in Im(f) \setminus \{\bot\}, \forall r \mbox{ s.t. } f(r) =  y, \forall i \in W: \\
        ~~\left | \left \{ r_i : \Pr[R_i = r_i|R_{<i} = r_{<i}, y]] \geq 0 \right \}  \right | \leq 2^{\delta} \cdot c.
    \end{eqnarray*}
    \item $\{\rho_i\}_{i \in V}$
    are distributed independently and uniformly at random
    conditioned on $f(R)$.
    % \mingyu{I still couldn't quite understand what $\rho_i$ represents, is it somehow to capture the fact that $R_i$ is not uniformly random?}
    % \dnote{It is capturing that the leakage functions can be probabilistic instead of deterministed. In particular, the leakage functions can depend on the output of a random oracle.}
\end{enumerate}
\end{theorem}
The proof is found in Section~\ref{sec:construct-gcp2}. 
%
Typically, one would like to set $c$ as large as possible, while ensuring that
the size of $V$ remains above some threshold. The achievable tradeoffs between $c$ and $|V|$ are determined by the min-entropy of $R$
before the spoiling bits $f(R)$ are released. For our applications, we require
$c = n^{1.5}$ and $|V| \geq n/3$. 
In Section~\ref{sec:tradeoff},
We show that our min-entropy assumption on $R$
implies that this parameter setting is achievable. 

